\documentclass[10pt]{article}
\usepackage[T1]{fontenc}
\usepackage[utf8]{inputenc}
\usepackage[french]{babel}
\usepackage{amsmath,amssymb,mathrsfs,mathpazo}
\usepackage{geometry}
\geometry{a4paper,margin=2cm}
\usepackage{graphicx,array,tabularx,tabularray,colortbl,multirow,multicol,booktabs,enumitem,pifont,ulem,fancybox,diagbox,cancel}
\usepackage{pstricks,pstricks-add,pst-plot,pst-tree,pst-node,pst-text,pst-func,pst-math,pst-bspline,pst-3dplot,pst-eucl,pst-solides3d}
\let\pscalculate\undefined
\usepackage{pst-bezier}
\usepackage{tikz,pgfplots,tkz-tab}
\usetikzlibrary{arrows,arrows.meta,patterns,calc,positioning,trees,decorations.markings,shapes.geometric,shapes.misc,shapes.arrows,intersections,decorations.pathmorphing,decorations.pathreplacing,backgrounds,shadows,babel}
\pgfplotsset{compat=1.18}
\usepgfplotslibrary{fillbetween}
\usepackage[np]{numprint}
\usepackage{listings,makecell,pas-tableur,twoopt}
\let\barre\undefined
\usepackage[pstricks]{tablvar}
\definecolor{darkgreen}{RGB}{0,100,0}
\providecommand{\red}{\color{red}}
\providecommand{\blue}{\color{blue}}
\providecommand{\green}{\color{green}}
\providecommand{\black}{\color{black}}
\providecommand{\white}{\color{white}}
\providecommand{\gray}{\color{gray}}
\providecommand{\cyan}{\color{cyan}}
\providecommand{\magenta}{\color{magenta}}
\providecommand{\yellow}{\color{yellow}}
\providecommand{\violet}{\color{violet}}
\providecommand{\orange}{\color{orange}}
\providecommand{\euro}{\texteuro}
\providecommand{\EUR}{\texteuro}
\providecommand{\ieme}{\textsuperscript{e}}
\providecommand{\degres}{\ensuremath{^\circ}}
\providecommand{\degree}{\ensuremath{^\circ}}
\providecommand{\vv}[1]{\overrightarrow{#1}}
\providecommand{\Coord}[1]{#1}
\providecommand{\decosix}{}
\providecommand{\starredbullet}{\ensuremath{\bullet}}
\pagestyle{empty}
\setlength{\parindent}{0pt}
\usepackage{adjustbox}\DeclareUnicodeCharacter{2236}{:}\DeclareUnicodeCharacter{2190}{\ensuremath{\leftarrow}}
\providecommand{\R}{\mathbb{R}}
\providecommand{\N}{\mathbb{N}}
\providecommand{\Z}{\mathbb{Z}}
\providecommand{\Q}{\mathbb{Q}}
\providecommand{\D}{\mathbb{D}}
\providecommand{\C}{\mathbb{C}}
\providecommand{\e}{{\mathrm{e}}}
\providecommand{\vect}[1]{\overrightarrow{#1}}
\providecommand{\vectt}[1]{\overrightarrow{\mathrm{#1}}}
\providecommand{\vvt}[1]{\overrightarrow{\mathrm{#1}}}
\providecommand{\np}[2][]{{#2~#1}}
\providecommand{\vv}[1]{{\overrightarrow{#1}}}
\renewcommand{\d}{{\mathrm{d}}}
\providecommand{\degres}{{^\circ}}
\providecommand{\degre}{{^\circ}}
\providecommand{\blue}{\color{blue}}
\providecommand{\red}{\color{red}}
\providecommand{\enskip}{\quad}
\providecommand{\llbracket}{\lbrack\!\lbrack}
\providecommand{\rrbracket}{\rbrack\!\rbrack}
\providecommand{\euro}{{\text{€}}}
\providecommand{\no}{{\text{n°}}}
\providecommand{\interval}[2]{{[#1;#2]}}
\providecommand{\intervFO}[2]{{[#1;#2[}}
\providecommand{\dotfill}{\ldots}

\providecommand{\ds}{\displaystyle}
\providecommand{\textding}[1]{#1}
\providecommand{\Oij}{\(\mathrm{O};\vec{\imath},\vec{\jmath}\)}
\providecommand{\Oijk}{\(\mathrm{O};\vec{\imath},\vec{\jmath},\vec{k}\)}
\providecommand{\Ouv}{\(\mathrm{O};\vec{u},\vec{v}\)}
\providecommand{\intervFF}[2]{\ensuremath{[#1;#2]}}
\renewcommand{\euro}{\ifmmode\text{\texteuro}\else\texteuro\fi}
\DeclareUnicodeCharacter{00B0}{\ifmmode^{\circ}\else\textdegree\fi}
\DeclareUnicodeCharacter{20AC}{\ifmmode\text{\texteuro}\else\texteuro\fi}
% Énoncés : Éducation nationale. Transcription et corrigés : APMEP, Denis Vergès et contributeurs. Découpage et classement : CoopMaths, Sébastien Lozano, Éric Elter et contributeurs. Adaptation de la présentation : Maths Explorer.
% CC BY-SA : https://coopmaths.fr/www/about/
% https://coopmaths.fr/alea/static/bac/2024/tex/bac_2024_06_sujet1_suede_2.tex
\begin{document}
\noindent\textit{APMEP / CoopMaths. CC BY-SA. Présentation adaptée par Maths Explorer.}\par\medskip
Au cours d'une séance, un joueur de volley-ball s'entraîne à faire des services. La probabilité qu'il réussisse le premier service est égale à 0,85.

On suppose de plus que les deux conditions suivantes sont réalisées:

\begin{itemize}[label=$\bullet$]
\item si le joueur réussit un service, alors la probabilité qu'il réussisse le suivant est égale à 0,6;
\item si le joueur ne réussit pas un service, alors la probabilité qu'il ne réussisse pas le suivant est égale à 0,6.
\end{itemize}

Pour tout entier naturel $n$ non nul, on note $R_{n}$ l'évènement \og le joueur réussit le $n$-ième service \fg{} et $\overline{R_{n}}$ l'évènement contraire.

\subsection*{Partie A}
On s'intéresse aux deux premiers services de l'entraînement.

\begin{enumerate}
	\item Représenter la situation par un arbre pondéré.
	\item Démontrer que la probabilité de l'évènement $R_{2}$ est égale à 0,57.
	\item Sachant que le joueur a réussi le deuxième service, calculer la probabilité qu'il ait raté le premier.
	\item Soit $Z$ la variable aléatoire égale au nombre de services réussis au cours des deux premiers services.
	\begin{enumerate}
		\item Déterminer la loi de probabilité de $Z$ (on pourra utiliser l'arbre pondéré de la question 1).
		\item Calculer l'espérance mathématique $\mathrm{E}(Z)$ de la variable aléatoire $Z$.

		Interpréter ce résultat dans le contexte de l'exercice.
	\end{enumerate}
\end{enumerate}

\subsection*{Partie B}

On s'intéresse maintenant au cas général.

Pour tout entier naturel $n$ non nul, on note $x_{n}$ la probabilité de l'évènement $R_{n}$.

\begin{enumerate}
	\item \begin{enumerate}
		\item Donner les probabilités conditionnelles $P_{R_{n}}\left(R_{n + 1}\right)$ et $P_{\overline{R_{n}}}\left(\overline{R_{n + 1}}\right)$.

		\item Montrer que, pour tout entier naturel non nul $n$, on a : $x_{n + 1}=0,2 x_{n} + 0,4$.
	\end{enumerate}

	\item Soit la suite $\left(u_{n}\right)$ définie pour tout entier naturel $n$ non nul par :\quad $u_{n}=x_{n}-0,5$.

	\begin{enumerate}
		\item Montrer que la suite $\left(u_{n}\right)$ est une suite géométrique.
		\item Déterminer l'expression de $x_{n}$ en fonction de $n$. En déduire la limite de la suite $\left(x_{n}\right)$.
		\item Interpréter cette limite dans le contexte de l'exercice.
	\end{enumerate}
\end{enumerate}

\end{document}
